% ============================================================
%  第六章
% ============================================================
\section{Sen 的试金石与 Seiberg--Witten 理论}
\markboth{第六章\ Sen 的试金石与 Seiberg--Witten 理论}{第六章\ Sen 的试金石与 Seiberg--Witten 理论}

\noindent{\small\color{shuyiGray}对应底本：原 p.29--p.36。}\par\medskip

\subsection{弦论里的对偶之潮}

\srcp{29}到 1990 年代初，围绕 Montonen--Olive 对偶的局面正在改变。对我来说，主要的改变是：同事们——Ashoke Sen、John Schwarz、Mike Duff、Paul Townsend、Chris Hull，以及不幸还有太多人，无法在这里一一列出——正在提出与 Montonen--Olive 对偶类似的猜想，但对象是弦论而不是量子场论。这些论证极具启发性。对我来说，当时还没有一锤定音的证据（\emph{smoking gun}），我也看不出人们怎能指望取得扎实的进展。但我已经不像 1977 年那么怀疑了。

\begin{yizhu}[smoking gun]
\textit{smoking gun} 直译“冒烟的枪”，出自“案发时枪还在冒烟”的典故，意为
“一锤定音、无可辩驳的证据”。原文用它形容 Sen 1994 年的文章；译文保留原文并括注，
后文再出现时径用“一锤定音的证据”。
\end{yizhu}


\subsection{1994 年 2 月 7 日}

\srcp{30}那件一锤定音的证据，于 1994 年 2 月 7 日提交到了 arXiv。

\begin{abstractbox}{Sen, \emph{Dyon--Monopole Bound States, Self-Dual Harmonic Forms on the Multi-Monopole Moduli Space, and $SL(2,\ZZ)$ Invariance in String Theory}（Physics Letters B329, 1994, pp.~217--221）}
Existence of $SL(2,\ZZ)$ duality in toroidally compactified heterotic string theory (or in the $\cN=4$ supersymmetric gauge theories), that includes the strong-weak coupling duality transformation, implies the existence of certain supersymmetric bound states of monopoles and dyons. We show that the existence of these bound states, in turn, requires the existence of certain normalizable, (anti-)self-dual, harmonic forms on the moduli space of BPS multi-monopole configurations, with specific symmetry properties. We give an explicit construction of this harmonic form on the two monopole moduli space, thereby proving the existence of all the required bound states in the two-monopole sector.

\medskip
环面紧致化的杂化弦理论（或 $\cN=4$ 超对称规范理论）中 $SL(2,\ZZ)$ 对偶的存在性——其中包括强—弱耦合的对偶变换——蕴含着单极子与双子（dyon）的某些超对称束缚态的存在。我们证明：这些束缚态的存在，反过来又要求 BPS 多单极子构型模空间上存在某些可归一化的、（反）自对偶的调和形式，并具有特定的对称性。我们在两单极子模空间上显式构造了这个调和形式，从而证明了两单极子部门里所有必需束缚态的存在。
\end{abstractbox}

\subsection{多单极子束缚态与 $L^2$ 调和形式}

\srcp{31}Sen 明白：$\cN=4$ 超杨--米尔斯理论的 Montonen--Olive 对偶，蕴涵着多单极子束缚态的存在。他证明，这样的束缚态对应于约化单极子模空间上中间维数的 $L^2$ 调和形式——那些 Atiyah 与 Hitchin 研究过的超 Kähler 流形。对两单极子的情形，Atiyah 与 Hitchin 已经显式描述了约化模空间，而他们的描述足够简单，使 Sen 能够显式地证明所需的 $L^2$ 调和形式确实存在。

\subsection{如遭电击}

\srcp{32}这篇论文像电击一样击中了我。一方面，我把它看作自早年（1977 年）以来对 Montonen--Olive 对偶的第一次根本性的新检验。在它通过这项检验之后，我真的不再认真怀疑它是否成立了。另一方面，Sen 表明了手头的工具真的足以取得实质进展。这激励我们所有人更加努力。

\subsection{与 Seiberg 的五个月}

\srcp{33}大约在此前一年，Nati Seiberg 在 $\cN=1$ 超杨--米尔斯理论的“动力学”（量子行为）上取得了重大进展。$\cN=1$ 超杨--米尔斯比 $\cN=2$ 更接近物理，但离 Montonen--Olive 对偶更远。Nati 邀我参与进来，尝试理解 $\cN=2$ 的量子行为。要记得，$\cN=2$ 正是对 Montonen--Olive 对偶“几乎”正确的那个理论。Sen 的论文以及某些定性考虑，把我们的注意力集中到电—磁对偶在理解 $\cN=2$ 动力学上可能扮演的角色。我们最终发现（在 Sen 的论文之后约五个月）的是：$\cN=2$ 由 Montonen--Olive 对偶的一种\emph{有效场论版本}所支配。由于希格斯机制，$\cN=2$ 理论在低能下可以一般地被一个有效的 $U(1)$ 理论描述；而正如我先前解释过的，$U(1)$ 规范理论确实有电—磁对偶。把这些陈述与某些定性考虑、以及关于 $\cN=2$ 超杨--米尔斯理论的一些已知事实结合起来，Seiberg 和我得到了这个理论的一个描述。

\subsection{场空间里的特殊点}

\srcp{34}从物理的观点看，我们工作的主要洞见是：尽管这个理论在低能下一般可以被描述为一个带质量电荷与单极子的自由 $U(1)$ 理论，但在场空间里有一些特殊的点，在那里单极子表现为轻的（甚至无质量的）基本粒子。这相当戏剧性，而且正合 Montonen 与 Olive 的精神：一个粒子——单极子——在弱耦合下被理解为延展对象、Bogomolny 方程的解，而在强耦合下可以表现得像一个点状的基本粒子。

\subsection{Seiberg--Witten 方程}

\srcp{35}把我们的结果应用到与 Donaldson 理论相关的 $\cN=2$ 理论的扭曲版本上，就意味着：Donaldson 不变量所包含的四维流形信息，也可以换一种方式得到——去数如今一般被称为 Seiberg--Witten 方程的解。Seiberg--Witten 方程是关于一个 $U(1)$ 联络 $A$ 与一个带电旋量场 $M$ 的方程：
\begin{equation}
  F^+=\overline{M}\,\Gamma M,\qquad \slashed{D}M=0 .
  \label{eq:sw}
\end{equation}
这里 $M$ 就是那个无质量单极子。在这个描述里，$M$ 看起来不像 Bogomolny 方程的解——而在 Atiyah 与 Hitchin 的书里，它正是 Bogomolny 方程的解。

\subsection{与 Langlands 的关系仍不明}

\srcp{36}此后又有许多发展；到 1990 年代中期，事情已经清楚：电—磁对偶不只是对的，而且在量子场论与弦论的一般理解中极为重要。然而，它与 Langlands 对应的关系仍然不清楚。

\begin{pizhu}{$SL(2,\ZZ)$ 与“必须存在的束缚态”}
$\cN=4$ 超杨--米尔斯理论有 $SL(2,\ZZ)$ 对偶群：$S$ 把 $\tau$ 送到 $-1/\tau$ 并交换电与磁，$T$ 把 $\tau$ 送到 $\tau+1$（即 $\theta\mapsto\theta+2\pi$）。在电荷格上，$SL(2,\ZZ)$ 以
\[
  \begin{pmatrix}m\\ n\end{pmatrix}\longmapsto
  \begin{pmatrix}a&b\\ c&d\end{pmatrix}\begin{pmatrix}m\\ n\end{pmatrix}
\]
的方式作用，其中 $(m,n)$ 是（磁荷，电电荷）。

弱耦合下有 $W$ 玻色子，电荷 $(0,1)$（在某个归一化下）。既然对偶是\emph{严格}的对称性，那么对 $W$ 玻色子做任意 $SL(2,\ZZ)$ 变换得到的态也必须存在。$S$ 把 $(0,1)$ 送到 $(1,0)$——这就是基本磁单极子；更一般地，$(0,1)$ 在 $SL(2,\ZZ)$ 下的轨道包含所有 $(m,n)$ 互素的态。这些态在弱耦合下不是基本场，而是“单极子加上 $n$ 个单位电电荷”的束缚态。它们的存在不是显然的：一个延展的单极子能否真的束缚住电荷，需要动力学上的答案。

Sen 的论证是：对偶蕴涵这些束缚态存在；而如果从模空间的角度看，这些束缚态就是单极子模空间上的 $L^2$ 调和形式。于是“对偶”这个关于耦合的断言，被翻译成了一个纯几何的、可以在弱耦合下独立检验的断言。这就是它为什么是一锤定音的证据：它不依赖对偶，却给出了对偶所预言的态。
\end{pizhu}

\begin{pizhu}{为什么是 $L^2$ 调和形式}
在 BPS 情形，一个多单极子构型的模空间 $\cM_n$ 上的零模（zero mode）参数化了它的变形。把量子力学限制在这个模空间上（即所谓“模空间量子力学”），波函数就是 $\cM_n$ 上的微分形式；超对称变换对应于外微分 $\dd$ 及其伴随 $\dd^*$。BPS 束缚态是超对称的基态，因此是 $\dd$ 与 $\dd^*$ 的共同零模——也就是\emph{调和形式}。又因为态必须是可归一化的，所以要求形式属于 $L^2$ 上同调。中间维数（$\cM_n$ 实维的一半）对应半超对称的态。

对两单极子，$\cM_2$ 是四维的 Atiyah--Hitchin 流形，中间维数是 $2$。Sen 要的是一族具有特定对称性的 $L^2$ 调和 $2$-形式；Atiyah--Hitchin 的显式度量使这件事可以逐项算出来。这就是“工具真的够用”的意思。
\end{pizhu}

\begin{pizhu}{Seiberg--Witten：$\cN=2$ 的精确解}
纯 $\cN=2$ 超杨--米尔斯（规范群 $SU(2)$）的低能有效理论是 $U(1)$ 的；它的 Coulomb 分支由一个复参数 $u$ 参数化（$u=\langle\operatorname{tr}\phi^2\rangle$ 的某种归一化）。Seiberg 与 Witten 发现，有效耦合 $\tau(u)$ 不是微扰论算出来的，而是某个椭圆曲线族的周期：
\[
  y^2=x^3-u\,x^2+\Lambda^4 x ,
\]
其中 $\Lambda$ 是强耦合标度。$\tau(u)$ 由这条曲线上的全纯微分 $\dd x/y$ 的周期比给出。曲线的判别式在 $u=\pm2\Lambda^2$ 处为零，对应两个奇异点：在一个点上，一个磁单极子变得无质量；在另一个点上，一个双子变得无质量。

“单极子凝聚”给出了夸克禁闭的模型：在磁单极子凝聚的真空里，电场被排斥到一根弦上（对偶迈斯纳效应），这正是 't Hooft 在 1970 年代设想的禁闭机制。这也是“强耦合下单极子可以像基本粒子”这句话的具体含义：在弱耦合下它是孤子，在强耦合下它是有效理论里的基本场。
\end{pizhu}

\begin{pizhu}{从 Donaldson 到 Seiberg--Witten}
Witten 在 1994 年的文章里证明：扭曲的 $\cN=2$ 理论在低能下退化为对 Seiberg--Witten 方程解的计数，于是 Donaldson 不变量可以由 Seiberg--Witten 不变量算出。这在数学上是一个巨大的简化：Donaldson 理论要数 ASD 联络模空间（维数高、奇点多），Seiberg--Witten 只要数 $U(1)$ 联络与旋量场的解（维数为零，本质上是有限个点）。

式 \eqref{eq:sw} 里，$F^+$ 是 $A$ 的自对偶部分，$M$ 是旋量（数学上属于某个 $\mathrm{Spin}^c$ 丛），$\Gamma$ 是 Clifford 乘法。这两条方程分别说“曲率由旋量的平方给出”与“旋量在联络下是和谐的”。它们是一阶椭圆系统，模空间是有限维的紧流形，因此不变量定义起来干净得多。

这一等价后来在数学上被证明（在适当的技术条件下，由 Feehan--Leness 等人完成），是这场物理—数学对话最直接的一次“兑现”。
\end{pizhu}
